% Lettre de motivation (FR) — El Mehdi Haress
% Cible : Analyste Quantitatif Risques de marché — périmètre Equity (réf. 26000I6J) — Société Générale
% Compilation : xelatex lettre-sg-quant.tex
\documentclass[10pt,a4paper]{article}

\usepackage{fontspec}
\usepackage[french]{babel}
\usepackage[T1]{fontenc}
\usepackage[
  top=1.5cm, bottom=1.3cm, left=2.0cm, right=2.0cm
]{geometry}
\usepackage{xcolor}
\usepackage{microtype}
\usepackage[hidelinks]{hyperref}
\usepackage{fontawesome5}

% ---------- Polices (identiques au CV) ----------
\setmainfont{texgyrepagella}[
  Extension      = .otf,
  Scale          = 1.05,
  UprightFont    = *-regular,
  BoldFont       = *-bold,
  ItalicFont     = *-italic,
  BoldItalicFont = *-bolditalic,
]
\newfontfamily\headingfont{texgyreheros}[
  Extension      = .otf,
  UprightFont    = *-regular,
  BoldFont       = *-bold,
  ItalicFont     = *-italic,
  BoldItalicFont = *-bolditalic,
]

% ---------- Couleurs (identiques au CV) ----------
\definecolor{accent}{RGB}{28,58,94}
\definecolor{ink}{RGB}{26,26,26}
\definecolor{muted}{RGB}{95,105,120}
\definecolor{rulegray}{RGB}{200,206,214}
\color{ink}

% Caractères laissés tels quels : un kern négatif casse l'extraction du mot-clé « C++ » par les ATS.
\newcommand{\cpp}{\mbox{C++}}

% ---------- Mise en page ----------
\pagestyle{empty}
\setlength{\parindent}{0pt}
\setlength{\parskip}{1.25ex plus 0.6ex minus 0.15ex}
\linespread{1.10}
\frenchspacing

\begin{document}

% =====================  EN-TÊTE  =====================
{\headingfont\bfseries\fontsize{20}{23}\selectfont EL MEHDI HARESS}\\[0.7ex]
{\headingfont\fontsize{10.5}{13}\selectfont\color{accent}%
Docteur en mathématiques appliquées \textbar{} Ingénieur CentraleSupélec}\\[1.0ex]
{\small
\faEnvelope[regular]~\href{mailto:el-mehdi.haress@inria.fr}{el-mehdi.haress@inria.fr}\quad
\faPhone*~+33 7 58 23 41 50\quad
\faLinkedin~\href{https://www.linkedin.com/in/el-mehdi-haress-67750b3aa/}{el-mehdi-haress}\quad
\faGlobe~\href{https://elmehdiharess.org}{elmehdiharess.org}
}

\vspace{0.4ex}
{\color{rulegray}\rule{\linewidth}{0.7pt}}

\vspace{1.2ex}

% =====================  DESTINATAIRE  =====================
\begin{minipage}[t]{0.5\textwidth}
\vspace{0pt}\raggedright\small
\textbf{Société Générale}\\
Validation des modèles \textemdash{} Opérations de marché\\
La Défense, Île-de-France
\end{minipage}\hfill
\begin{minipage}[t]{0.42\textwidth}
\vspace{0pt}\raggedleft\small
Sophia Antipolis, le 15 septembre 2026
\end{minipage}

\vspace{2.2ex}

{\small\textbf{Objet :} candidature au poste d'Analyste Quantitatif Risques de marché \textemdash{} périmètre
Equity (référence 26000I6J)}

\vspace{2.2ex}

Madame, Monsieur,

\og~Vous aimez vous plonger dans les équations mathématiques […] mais vous n'aimez pas l'atmosphère studieuse
des bibliothèques ?~\fg{} J'ai passé cinq ans dans ces bibliothèques, et j'y ai démontré des théorèmes dont je
suis fier. C'est précisément parce que j'aime ces objets que je souhaite aujourd'hui les voir confrontés à des
books, à des positions et à des décisions qui engagent \textemdash{} plutôt qu'à un referee anonyme.

Cette bascule est moins un changement de métier qu'un changement de terrain. Mes recherches portent sur les
\textbf{équations différentielles stochastiques dirigées par un mouvement brownien fractionnaire}, c'est-à-dire
sur des dynamiques à mémoire longue et à trajectoires irrégulières. Ce sont exactement les objets sur lesquels
repose la \textbf{volatilité rough} \textemdash{} rough Heston, rough Bergomi \textemdash{} dont la littérature
est née de l'observation des surfaces de volatilité actions. Mieux : l'un de mes articles construit des
estimateurs conjoints de la dérive, de la volatilité et de \textbf{l'indice de Hurst} sur observations
discrètes, avec vitesses de convergence prouvées. Autrement dit, j'ai travaillé sur la question que pose un
validateur devant une calibration de volatilité : \emph{cet estimateur converge-t-il, à quelle vitesse, et sous
quelles hypothèses se dégrade-t-il ?}

Ma thèse prolonge ce réflexe. Elle établit l'un des premiers cadres rigoureux d'approximation numérique pour
des dynamiques \textbf{à dérive singulière}, là où les hypothèses habituelles des schémas d'Euler tombent, et
mes travaux à Leeds quantifient la \textbf{sensibilité d'une solution aux perturbations de sa spécification et
de ses paramètres}. C'est, dans un langage académique, le cœur du risque de modèle : identifier où un modèle
cesse d'être fiable et le démontrer. J'ai par ailleurs abordé le risque par l'autre bout lors de mon séjour à
l'Université d'Osaka, en transposant les \textbf{mesures de risque spectrales} (VaR, CVaR) hors de leur cadre
financier d'origine, avec des garanties sous queues lourdes et événements rares.

Enfin, votre annonce insiste sur l'oral et l'écrit, et c'est un terrain que je connais. J'ai enseigné la
\textbf{finance stochastique et la modélisation du risque en M2 à CentraleSupélec}, après avoir donné plus de
quinze exposés internationaux et co-organisé deux rencontres scientifiques internationales. Expliquer un
modèle, expliciter ses hypothèses et
assumer ses limites devant un auditoire qui va vous contredire : c'est l'exercice de l'enseignement, et c'est
aussi, si je comprends bien, celui d'un rapport de validation défendu devant des traders et des quants.

Ce qui m'attire dans votre offre est précisément le \textbf{quatuor d'analyses} que vous décrivez. Valider un
modèle de valorisation, donner un go sur un nouveau produit, challenger une méthode d'estimation de paramètres
et concevoir des détecteurs de risque de modèle : ce sont les quatre questions qui structurent mon travail
depuis ma thèse, mais posées sur des positions réelles, avec une contradiction immédiate et une échéance sur
laquelle on ne négocie pas. Côté outils, je
développe en Python quotidiennement depuis huit ans, j'ai les bases du \cpp{} de ma formation d'ingénieur, et
j'ai construit seul plusieurs applications web de bout en bout, jusqu'à leur mise en production pour un client
réel : monter en compétence sur le C\# et contribuer à votre librairie de calcul interne ne me fait pas peur,
au contraire.

Je serais très heureux de vous exposer de vive voix ce que la rigueur du singulier et du fractionnaire peut
apporter à la détection du risque de modèle sur un périmètre Equity.

Veuillez agréer, Madame, Monsieur, l'expression de mes salutations distinguées.

\vspace{2.6ex}

\hfill
\begin{minipage}{0.42\textwidth}
\raggedleft
{\headingfont\bfseries\small El Mehdi Haress}\\[0.2ex]
{\small\color{muted}Docteur en mathématiques appliquées}
\end{minipage}

\end{document}
